%% NIM A style draft — Nuclear Instruments and Methods in Physics Research A
%% Compile with: pdflatex paper_draft.tex  (requires elsarticle.cls)
%% This is a WORKING DRAFT. Sections marked (TO BE DONE) are pending.

\documentclass[preprint,12pt]{elsarticle}

\usepackage{amsmath,amssymb}
\usepackage{graphicx}
\usepackage{booktabs}
\usepackage{siunitx}
\usepackage{xcolor}
\newcommand{\todo}[1]{{\color{red}\textbf{(TO BE DONE: #1)}}}

\journal{Nuclear Instruments and Methods in Physics Research A}

\begin{document}

\begin{frontmatter}

\title{Single-orbit estimation of the dominant false-EDM harmonic
in a frozen-spin proton storage ring}

\author[inst1]{S.~Hac\i\"omero\u{g}lu\corref{cor1}}
\cortext[cor1]{Corresponding author}
\ead{selcuk.haciomeroglu@gmail.com}

\address[inst1]{\todo{affiliation}}

\begin{abstract}
The proton electric dipole moment (pEDM) frozen-spin experiment requires
the vertical position of every storage-ring quadrupole to be known to
$\sim\SI{10}{\micro\meter}$, because uncorrected misalignments generate a
closed-orbit distortion that mimics a true EDM signal. Beam-position
monitors (BPMs) carry static electronic offsets of $\sim\SI{100}{\micro
\meter}$ that bury the alignment signal. The established mitigation,
$K$-modulation (kmod), differences two lattice configurations to cancel
the offset but pays a steep rank/conditioning penalty when only a few
quadrupoles can be modulated. We show that, for the specific goal of
estimating the dominant false-EDM-driving harmonic of the misalignment
pattern ($k=2$ for our lattice)---without attempting a full reconstruction
of the 48-dimensional misalignment vector---a single closed-orbit
measurement analysed with a response-matrix Fourier basis recovers the
target harmonic to $0.6\%$ in amplitude and \SI{0.046}{\radian} in phase
under a \SI{100}{\micro\meter} BPM offset, with no gradient modulation.
We trace this to three quantitative properties of the response matrix
$R$: (i) tune-proximity resonant amplification of the $k=2$ mode
($\|RF_{k=2}\|$ a factor $\sim$34 larger than the lattice average),
(ii) the consequent inversion of the misalignment-amplitude hierarchy at
the orbit level, and (iii) near-orthogonality of the orbit responses of
distinct harmonics ($|\langle RF_2,RF_{4,6,8}\rangle|\approx0.01$), which
prevents leakage of the large contaminant harmonics into the $k=2$
estimate. We also report a unit-handling error in the response-matrix
pipeline that had inflated $R$ by $10^3$ and its correction.
\end{abstract}

\begin{keyword}
storage ring \sep electric dipole moment \sep closed-orbit distortion \sep
quadrupole alignment \sep response matrix \sep beam position monitor
\end{keyword}

\end{frontmatter}

%==============================================================
\section{Introduction}
\label{sec:intro}

A storage-ring search for the proton electric dipole moment (pEDM)
operates in the frozen-spin regime, in which the in-plane spin precession
relative to the momentum is cancelled so that a non-zero EDM produces a
slow, accumulating vertical spin rotation. The dominant systematic
limitation is the vertical closed-orbit distortion (COD): a quadrupole
displaced vertically by $\Delta y_j$ deflects the beam, and the resulting
COD samples radial fields that rotate the spin in the same way a true EDM
would. Quantitatively, suppressing this false signal below the target
sensitivity requires knowledge of each quadrupole's vertical position at
the $\sim\SI{10}{\micro\meter}$ level, well beyond mechanical survey
precision; a beam-based determination is therefore mandatory.

The measurement model is linear,
\begin{equation}
\mathbf{y} = R\,\Delta q + \mathbf{b} + \boldsymbol{\eta},
\label{eq:model}
\end{equation}
where $\mathbf{y}\in\mathbb{R}^{48}$ are the BPM readings, $R$ is the
$48\times48$ orbit response matrix, $\Delta q$ the unknown quadrupole
misalignments, $\mathbf{b}$ the static BPM electronic offsets
($\sim\SI{100}{\micro\meter}$, constant over hours--days but unknown), and
$\boldsymbol{\eta}\sim\SI{1}{\micro\meter}$ the readout noise. The
difficulty is entirely $\mathbf{b}$: a direct inversion
$\Delta q = R^{-1}(\mathbf{y}-\mathbf{b})$ requires knowing $\mathbf{b}$,
which is unavailable; solving for $\Delta q$ alone from $R^{-1}\mathbf{y}$
carries an $R^{-1}\mathbf{b}$ contamination that, with $\kappa(R)\approx249$,
reaches the millimetre level---two orders of magnitude above the target.

The conventional remedy is $K$-modulation (kmod): the lattice is measured
in two gradient settings and the readings differenced,
\begin{equation}
\Delta\mathbf{y} = \mathbf{y}_2-\mathbf{y}_1
= (R_2-R_1)\,\Delta q + (\boldsymbol{\eta}_2-\boldsymbol{\eta}_1)
\equiv \Delta R\,\Delta q + \text{noise},
\end{equation}
which removes $\mathbf{b}$ exactly. When all 48 quadrupoles are modulated
uniformly, $\Delta R=\varepsilon R$ and $\kappa(\Delta R)=\kappa(R)$;
the method works but is operationally hard. When only one or two
quadrupoles can be modulated, $\Delta R$ becomes rank-deficient
($\mathrm{rank}(\delta R)\le$ number of modulated quadrupoles, since to
first order $\delta R = -R\,\mathrm{diag}(\delta K)\,R$), the condition
number rises to $\sim10^6$, and the reconstruction degrades severely.

In this work we step back from kmod and ask a narrower, physically
motivated question. \textbf{The goal is not a full reconstruction of the
48-dimensional misalignment vector $\Delta q$, but rather to estimate the
single dominant false-EDM-driving harmonic ($k=2$), which can subsequently
be compensated by dedicated correction elements (e.g.\ harmonic dipole
correctors distributed around the ring).} The design of that correction
system is a separate problem and lies outside the scope of this paper; here
we address only the measurement of the harmonic to be corrected. The
underlying premise is operational rather than metrological: \emph{a full
alignment reconstruction is unnecessary when the objective is suppression of
the dominant false-EDM harmonic.} This inverts the conventional
``measure-all-then-correct'' sequence---instead of resolving every
misalignment and then correcting, we measure directly the one mode that
matters for the systematic and target it.

The false EDM is not driven
equally by all 48 degrees of freedom of $\Delta q$; it is dominated by the
low-order Fourier harmonics of the misalignment pattern, and for our lattice
the leading contribution is the $k=2$ harmonic---a fact we confirm directly
by spin tracking in Sec.~\ref{sec:falseedm}, qualitatively consistent with
the $N=2$ peak of Omarov \emph{et al.}~\cite{omarov}. We therefore target
$k=2$ directly and show that it can be extracted from a \emph{single} orbit
measurement governed by Eq.~(\ref{eq:model}), keeping $\mathbf{b}$ in the
model rather than cancelling it. The key realisation is that the fitting-matrix column norm $\|M_{k=2}\|=167$
is 10.6$\times$ larger than $\|M_{k=4}\|=16$, so the $k=2$ estimate is
inherently less contaminated by any $\mathbf{b}$---regardless of its
statistical structure---than the $k=4$ estimate.
For any fixed $\mathbf{b}$, the $k=2$ error is
$\delta a_{k=2}=(\mathbf{b}\cdot\hat{m}_{k=2})/\|M_{k=2}\|$;
the safe limit is $\|\mathbf{b}_{\parallel}\|<\SI{1669}{\micro\meter}$,
independent of any whiteness assumption.
Section~\ref{sec:sim} describes the simulation,
Section~\ref{sec:method} the reconstruction, Section~\ref{sec:results}
the quantitative mechanism and results, and Section~\ref{sec:future} the
open items.

%==============================================================
\section{Simulation framework}
\label{sec:sim}

The ring consists of 24 FODO cells ($R_0=\SI{95.49}{\meter}$), each
holding one focusing (QF) and one defocusing (QD) quadrupole of length
\SI{0.4}{\meter} separated by \SI{2.083}{\meter} drifts and electrostatic
bending sections, for a total of 48 quadrupoles and 48 co-located BPMs.
The vertical and horizontal tunes are $Q_{y}\approx Q_{x}\approx2.68$.

Particle tracking uses a fourth-order Gauss--Legendre (GL4) symplectic
integrator implemented in C++ (\texttt{integrator.cpp}) and driven from
Python via \texttt{ctypes}. The symplectic scheme conserves phase-space
volume to machine precision, which is essential for the long multi-turn
runs from which the closed orbit is extracted. The response matrix is
built column by column: each quadrupole is displaced by
$\delta=\SI{e-4}{\meter}$, the closed orbit at the 48 BPMs is recorded,
and
\begin{equation}
R[:,j] = \frac{\mathbf{y}_\mathrm{CO}(\Delta y_j=\delta)
              -\mathbf{y}_\mathrm{CO}(0)}{\delta}.
\end{equation}

\paragraph{Unit-handling correction.}
During this study we identified a unit error in the response-matrix
pipeline. The tracker writes the orbit in millimetres
(\texttt{integrator.cpp} emits columns labelled \texttt{x\_mm},
\texttt{y\_mm}), whereas the reader \texttt{read\_cod\_quads} treated the
values as metres, so that $R=\text{orbit[mm]}/\text{misalign[m]}$ was
inflated by a factor $10^3$ ($\max|R|\approx1950$,
$\sigma_{\max}\approx34729$). Because the condition number and rank are
scale-invariant, kmod results---in which the factor cancels between
$\Delta\mathbf{y}$ and $\Delta R$---were unaffected; however, any test
mixing the physical BPM offset $\mathbf{b}$ with $R\,\Delta q$ was off by
$10^3$ in relative scale. After correcting the reader (a single
mm$\rightarrow$m conversion), $\max|R|\approx1.95$,
$\sigma_{\max}\approx34.73$, $\sigma_{\min}\approx0.140$, and
$\kappa(R)=248.7$; all matrices were rebuilt and the results below use the
corrected $R$.

\paragraph{Physical structure of the response matrix.}
It is instructive to decompose $R$ as
\begin{equation}
  R = G\,\mathrm{diag}(K_j),
  \label{eq:R_decomp}
\end{equation}
where $\mathrm{diag}(K_j)$ is the diagonal \emph{displacement-to-kick} gain
matrix ($K_j = k_j\ell_j$, the product of normalised gradient and effective
length) and $G$ is the betatron \emph{Green function} that converts a thin
vertical kick applied at element $j$ into the resulting closed-orbit
displacement observed at BPM $i$:
\begin{equation}
  G_{ij}
  = \frac{\sqrt{\beta_i\,\beta_j}}{2\sin\pi Q_y}
    \cos\!\bigl(|\mu_i-\mu_j|-\pi Q_y\bigr).
  \label{eq:green}
\end{equation}
Here $\beta_i$, $\mu_i$ are the Courant--Snyder amplitude and betatron
phase advance at BPM $i$, and $Q_y\approx2.68$ is the vertical tune.
Equation~(\ref{eq:green}) follows from the standard one-turn-map solution
of the Hill equation~\cite{lee_accel_phys}; it reveals that $G$ is
\emph{nearly symmetric} ($G_{ij}\approx G_{ji}$), with a numerically
verified relative asymmetry of $0.005$.
One can therefore read $R$ as the \emph{gain function of magnetic forces
along the lattice onto the BPMs}: the first factor $\mathrm{diag}(K_j)$
translates a physical displacement into an angular kick, while $G$
propagates that kick through the optics to produce an orbit response.

The decomposition~(\ref{eq:R_decomp}) makes two key facts transparent.
\begin{enumerate}
\item \emph{Why $\mathbf{b}$ bypasses $R$.}
  The BPM electronic offset $\mathbf{b}$ is added in measurement space,
  \emph{after} the mechanical kick chain
  $\Delta q \xrightarrow{\;\mathrm{diag}(K_j)\;} \boldsymbol{\theta}
             \xrightarrow{\;G\;} \mathbf{y}_\mathrm{orbit}$.
  It is never acted upon by $G$; it enters the observation
  $\mathbf{y}=G\,\mathrm{diag}(K_j)\,\Delta q+\mathbf{b}$ already in the
  output frame and therefore receives no lattice amplification.

\item \emph{Why $k=2$ is resonant.}
  The QF/QD sign alternation encoded in $K_j$ (and in the $(-1)^j$ factor of
  $F_k$) means that a smooth $k$-harmonic misalignment drives $G$ with a kick
  pattern at spatial frequency content near the vertical tune $Q_y$.
  The Green function $G$ amplifies most strongly the kick component whose
  frequency matches $Q_y$; because $Q_y\approx2.68$, the $k=2$ FODO harmonic
  is the closest resonant match.
  This is why $\|M_{k=2}\|=167$ while $\|M_{k=4}\|=16$ — a factor-of-ten
  difference that directly governs the BPM-offset robustness hierarchy derived
  in Section~\ref{sec:whiteness}.
\end{enumerate}

%==============================================================
\section{Method: Fourier-response projection}
\label{sec:method}

We parametrise the misalignment by a FODO-aware Fourier basis. Because the
QF and QD quadrupoles within a cell contribute COD with opposite sign,
the basis vectors carry an alternating factor,
\begin{equation}
F_k[j] = (-1)^{j}\cos\!\left(\frac{2\pi k\lfloor j/2\rfloor}{24}\right),
\end{equation}
and similarly for the sine component. Writing $\Delta q = F\hat{a}$,
Eq.~(\ref{eq:model}) becomes
\begin{equation}
\mathbf{y} = \underbrace{R F}_{M}\,\hat{a} + \mathbf{b} + \boldsymbol{\eta}.
\end{equation}
Rather than cancelling $\mathbf{b}$, we fit $\hat{a}$ directly. The core
operation is a projection of the orbit onto the columns of $M=RF$: because
the $k=2$ column norm $\|M_{k=2}\|=167$ is resonantly amplified (Section~\ref{sec:results}),
the projection is robust to the BPM offset $\mathbf{b}$ regardless of its
statistical structure, as shown by Eq.~(\ref{eq:bias}). In the targeted
scenario where $k=2$ is known in advance, a single weighted projection
$\hat{a}_{k=2}=M_{k=2}^T\mathbf{y}/\|M_{k=2}\|^2$ is sufficient.

For oracle-free operation---where no prior knowledge of which harmonics are
present is assumed---we employ a CLEAN iteration borrowed from radio
astronomy~\cite{clean}. With residual $\mathbf{r}\leftarrow\mathbf{y}$ and
a candidate set $k\in\{1,\dots,12\}$, each iteration projects $\mathbf{r}$
onto every candidate column of $M$, selects the largest, subtracts a
fraction $g=0.2$ of the match, and accumulates the amplitude.
CLEAN and direct least squares are statistically identical for the targeted
$k=2$ estimate ($k=2$ error ratio 1.0$\times$; $N=200$ Monte-Carlo):
the CLEAN wrapping adds oracle-free candidate selection at no cost to accuracy.

%==============================================================
\section{Results}
\label{sec:results}

\subsection{Direct confirmation that $k=2$ dominates the false EDM}
\label{sec:falseedm}

The premise of this work---that $k=2$ is the harmonic worth resolving---rests
on a physical claim: among equal-amplitude misalignment harmonics, $k=2$
produces the largest false EDM. We verify this directly by spin tracking,
independently of any reconstruction. For each $k=0,\dots,12$ we set
$\Delta y_j=A\,F_k[j]$ at equal amplitude $A=\SI{10}{\micro\meter}$, switch
off the true EDM, and measure the secular vertical-spin precession rate
$dS_y/dt$, which is the false-EDM observable.

Resolving $dS_y/dt$ is delicate: the secular drift ($\sim\!10^{-9}$\,rad/s)
is buried under an in-turn spin oscillation of order $10^{-5}$ in $S_y$.
A continuous Savitzky--Golay slope is then sampling-dependent and even
sign-unstable, so we remove the contamination at its source.
\emph{(i)}~The particle is launched on the distorted closed orbit, located
by a single Newton step that minimises the fixed-azimuth betatron variance
(an exact quadratic form for the linear lattice); this drives the residual
betatron amplitude to $\sim\!10^{-14}$\,m, so the orbit---and hence the
in-turn spin precession---is exactly periodic. \emph{(ii)}~$S_y$ is then
sampled \emph{stroboscopically}, once per turn at a fixed azimuth, where the
periodic part cancels identically and only the secular drift survives. The
rate follows from a straight-line fit, with no smoothing and no
sampling-rate choice.

Table~\ref{tab:falseedm} shows the outcome. The false EDM peaks sharply at
$k=2$, exceeding $k=1$ by $2.9\times$ and $k=3$ by $2.2\times$, and tracks
the closed-orbit amplitude (both maximal at $k=2$, the harmonic nearest the
vertical tune $Q_y\approx2.68$). The rate falls steeply for $k\ge3$,
reaching $4\%$ of the $k=2$ peak at $k=6{,}7$ and levelling at
$\lesssim\!2\%$ for $k\ge9$, where the CO amplitude is below
\SI{3}{\micro\meter}. This is qualitatively consistent with the
$N=2$ peak reported by Omarov \emph{et al.}~\cite{omarov} (their
Figs.~7--8), though the two rings differ in lattice design and the
comparison is not direct. There the low-harmonic rates were quoted only as
statistical upper limits, because the exact precession rate could not be
determined from the simulation; this is precisely the
oscillation-contamination problem that the closed-orbit\,/\,stroboscopic
method removes here, returning the actual rates.

\begin{table}[t]
\centering
\caption{False-EDM secular rate $dS_y/dt$ from a pure single-harmonic
quadrupole misalignment ($A=\SI{10}{\micro\meter}$, true EDM off), measured
by stroboscopic spin tracking on the closed orbit ($k=0$--$12$,
$t_2=\SI{2}{\milli\second}$, 20-turn CO correction). The closed-orbit
amplitude and the false EDM both peak at $k=2$.}
\label{tab:falseedm}
\begin{tabular}{cccc}
\toprule
$k$ & CO amplitude [mm] & $dS_y/dt$ [rad/s] & $|dS_y/dt|/|dS_y/dt|_{k=2}$ \\
\midrule
0  & 0.058 & $+1.80\times10^{-10}$ & 0.127 \\
1  & 0.070 & $+4.88\times10^{-10}$ & 0.344 \\
\textbf{2}  & \textbf{0.198} & $\mathbf{+1.42\times10^{-9}}$  & \textbf{1.000} \\
3  & 0.088 & $-6.40\times10^{-10}$ & 0.452 \\
4  & 0.028 & $-2.09\times10^{-10}$ & 0.147 \\
5  & 0.014 & $-9.22\times10^{-11}$ & 0.065 \\
6  & 0.008 & $-5.71\times10^{-11}$ & 0.040 \\
7  & 0.005 & $-5.26\times10^{-11}$ & 0.037 \\
8  & 0.004 & $-3.86\times10^{-11}$ & 0.027 \\
9  & 0.003 & $-2.26\times10^{-11}$ & 0.016 \\
10 & 0.002 & $-2.49\times10^{-11}$ & 0.018 \\
11 & 0.002 & $-2.24\times10^{-11}$ & 0.016 \\
12 & 0.002 & $-2.40\times10^{-11}$ & 0.017 \\
\bottomrule
\end{tabular}
\end{table}

The measurement is numerically robust. A null run ($A=0$) returns $dS_y/dt$
identically zero, so there is no spurious background. The response is exactly
odd under $\Delta y\to-\Delta y$ (ratio $-1.000$), confirming a clean
first-order coupling; the spin norm is conserved to $10^{-12}$ over the run;
the result reproduces across independent machines to better than $1\%$; and
halving the integration step ($10^{-11}\!\to\!5\times10^{-12}$\,s) shifts the
$k=2$ rate by only $2.7\%$ while leaving the $k_2/k_n$ ratios---and the
$k=2$ peak---unchanged. The closed-orbit-launch and stroboscopic procedure
is documented in \texttt{false\_edm\_yontemi.md}.

\subsection{Reconstruction test case}

We test against a structured truth: $k=2$ at \SI{10}{\micro\meter}
(the target) accompanied by large contaminants $k=4,6,8$ at
\SIrange{200}{300}{\micro\meter}, all with FODO antisymmetry, plus an
independent Gaussian BPM offset of $\sigma_b=\SI{100}{\micro\meter}$ per
monitor. The CLEAN candidate set is $k=1,\dots,12$. Results are averaged
over 50 Monte-Carlo realisations of $\mathbf{b}$.

\subsection{Why a weak $k=2$ survives among large high-$k$ harmonics}

The reconstruction succeeds because of three quantitative properties of
$R$, all directly computable.

\paragraph{(1) Tune-proximity resonant amplification.}
Each unit-norm Fourier mode $F_k/\|F_k\|$ is mapped to an orbit of norm
$\|RF_k\|$ (Table~\ref{tab:gain}). For the fitting matrix $M=RF$, the
relevant quantity is the column norm $\|M_k\|=\sqrt{24}\,\|RF_k\|$.
The $k=2$ column norm $\|M_{k=2}\|=167$ exceeds $\|M_{k=4}\|=16$
by a factor 10.6, and $\|M_{k=8}\|=2.9$ by a factor 58.
This directly controls contamination tolerance: a given offset $\mathbf{b}$
biases the $k=2$ estimate 10.6$\times$ less than the $k=4$ estimate.

\begin{table}[t]
\centering
\caption{Orbit gain $\|RF_k\|$ for unit-norm Fourier modes, and column norms
$\|M_k\|=\sqrt{24}\,\|RF_k\|$ of the fitting matrix $M=RF$.
The tune $Q_y\approx2.68$ places $k=2,3$ closest to resonance.
The column norm governs offset contamination: $|\delta a_k|\le\|\mathbf{b}_{\parallel}\|/\|M_k\|$.}
\label{tab:gain}
\begin{tabular}{ccccccc}
\toprule
$k$ & 1 & \textbf{2} & 3 & 4 & 6 & 8 \\
\midrule
$\|RF_k\|$ &  8.8 & \textbf{34.1} &  8.9 & 3.2 & 1.1 & 0.59 \\
$\|M_k\|$  & 43.0 & \textbf{167}  & 43.6 &  16 & 5.5 & 2.9  \\
\bottomrule
\end{tabular}
\end{table}

\paragraph{(2) Inversion of the amplitude hierarchy at the orbit level.}
Although the contaminants are $20$--$30\times$ larger in misalignment,
the amplification gap narrows or reverses the hierarchy once mapped to
orbit space (Table~\ref{tab:orbit}). The \SI{10}{\micro\meter} $k=2$
misalignment produces a \SI{1669}{\micro\meter} orbit norm---this is simply
$\|M_{k=2}\|\times a_{k=2}=167\times\SI{10}{\micro\meter}$, i.e.\ the column
norm times the misalignment amplitude. The orbit norm is comparable to $k=6$
and larger than $k=8$, so $k=2$ is no longer a buried weak signal but one
of the leading orbit components.

\begin{table}[t]
\centering
\caption{Misalignment amplitude versus orbit-norm contribution for the
truth harmonics.}
\label{tab:orbit}
\begin{tabular}{cccc}
\toprule
$k$ & misalignment & orbit norm & gain \\
\midrule
2 & \SI{10}{\micro\meter}  & \SI{1669}{\micro\meter} & 34.1 \\
4 & \SI{300}{\micro\meter} & \SI{4707}{\micro\meter} & 3.2 \\
6 & \SI{300}{\micro\meter} & \SI{1651}{\micro\meter} & 1.1 \\
8 & \SI{200}{\micro\meter} & \SI{577}{\micro\meter}  & 0.59 \\
\bottomrule
\end{tabular}
\end{table}

\paragraph{(3) Near-orthogonality of harmonic orbit responses.}
The normalised orbit patterns of distinct harmonics are almost
orthogonal,
\begin{equation}
|\langle RF_2, RF_{4,6,8}\rangle| \approx 0.01,
\end{equation}
so least squares separates $k=2$ from the contaminants cleanly: the large
$k=4,6,8$ amplitudes do not leak into the $k=2$ estimate. This is the
converse of the kmod situation, where the small-difference matrix
$\Delta R$ mixes the modes.

Together, these three properties ensure that the projection of $\mathbf{b}$
onto the $k=2$ column of $M$ is small relative to the signal---regardless
of the statistical structure of $\mathbf{b}$ (see Section~\ref{sec:whiteness})---and yield the reconstruction below.

\subsection{A matched-filter / conditioning view}
\label{sec:matchedfilter}

The three properties above have a compact linear-algebra interpretation that
explains why the same resonance that makes $k=2$ the dominant false-EDM driver
also makes it the most offset-robust harmonic to estimate. The single-mode
estimator is the Moore--Penrose pseudoinverse of one column,
\begin{equation}
\hat a_{k=2} = \frac{M_{k=2}^T\,\mathbf{y}}{\|M_{k=2}\|^2}
             = M_{k=2}^{+}\,\mathbf{y},
\qquad
M_{k=2}^{+} = \frac{M_{k=2}^T}{\|M_{k=2}\|^2},
\end{equation}
and the operator norm of the pseudoinverse of a vector is the reciprocal of
its norm, $\|M_{k=2}^{+}\|=1/\|M_{k=2}\|=1/167$.

\paragraph{Reciprocity.} Signal and offset enter the data at different points,
so they experience different net gains. A genuine $k=2$ misalignment
$\Delta q=A\,F_{k=2}$ propagates \emph{forward} through $R$ (gain
$\|M_{k=2}\|$) and is then back-projected through $M_{k=2}^{+}$ (gain
$1/\|M_{k=2}\|$); the two cancel exactly, $M_{k=2}^{+}(A\,M_{k=2})=A$. The BPM
offset $\mathbf{b}$ never passes forward through $R$---it is added directly in
measurement space---so it experiences \emph{only} the back-projection:
$|M_{k=2}^{+}\mathbf{b}|\le\|\mathbf{b}\|/\|M_{k=2}\|$. Whatever traverses $R$
both ways is recovered at unit gain; whatever enters only on the measurement
side is suppressed by $\|M_{k=2}\|$.

\paragraph{Conditioning.} In SVD terms $R=U\Sigma V^T$ with
$\sigma_{\max}=34.7$, $\sigma_{\min}=0.14$, $\kappa(R)=249$. Crucially, the
two leading right singular vectors of $R$ are almost exactly the $k=2$ cosine
and sine Fourier modes (verified by projecting each $V[:,i]$ onto all $F_k$,
$k=0,\ldots,12$, with $>99.5\%$ of norm captured at $k=2$), confirming that
$k=2$ sits at the top of the singular spectrum because it is the mode closest
to the tune $Q_y\approx2.68$ resonance. The $k=2$
misalignment direction is almost entirely aligned with the leading
right-singular vector ($91\%$ of $\|F_{k=2}\|^2$, and $\|RF_{k=2}\|=34.1\approx
\sigma_{\max}$): $k=2$ sits at the \emph{top} of the spectrum of $R$ because it
is the mode closest to the $Q_y\approx2.68$ resonance. A naive full inversion
$R^{-1}=V\Sigma^{-1}U^T$ is dangerous precisely because it is dominated by
$1/\sigma_{\min}$ and amplifies any offset by up to $\kappa=249$ (the
millimetre-level contamination of Section~\ref{sec:intro}). Targeting a single
resonant mode replaces that ill-conditioned $48$-dimensional inversion with a
one-dimensional projection of condition number $1$, operating in the
large-singular-value subspace where the pseudoinverse \emph{contracts} rather
than amplifies the offset.

\paragraph{Matched-filter gain.} Equivalently, the estimator is a matched
filter whose template is the $k=2$ orbit fingerprint $M_{k=2}$; its
signal-to-noise gain is $\|M_{k=2}\|$, giving
\begin{equation}
\mathrm{SNR}_{\mathrm{out}}
= \frac{A\,\|M_{k=2}\|}{\sigma_b}
= \frac{\SI{10}{\micro\meter}\times167}{\SI{100}{\micro\meter}}
\approx 16.7 .
\end{equation}
The resonant amplification thus works doubly in the estimator's favour: it
enlarges the genuine $k=2$ signal in BPM space ($A\|M_{k=2}\|=\SI{1670}{\micro
\meter}$) while simultaneously dividing the offset contribution by the same
$\|M_{k=2}\|=167$.

\subsection{Coverage of external radial field harmonics}
\label{sec:bxequiv}

The method targets the $k$-th harmonic of the \emph{effective} closed-orbit
distortion irrespective of its physical origin.  A ring-wide radial magnetic
field $B_x(\theta)=B_r\cos(k\theta)$ drives the same Fourier component of
the orbit distortion as a $k$-th harmonic quad misalignment, because both
couple to the spin through the same ring-integrated radial impulse.  The
framework therefore covers external field harmonics (systematic-field
imperfections, residual stray fields) on equal footing with mechanical
misalignments: a single-orbit measurement of the $k=2$ harmonic amplitude
suppresses the false EDM regardless of whether the distortion is mechanical
or magnetic in origin.

We express sensitivity results in terms of the spin-equivalent radial
field $B_{x,\mathrm{eq}}^{(k)}$, the amplitude of $B_x(\theta)=B_{x,\mathrm{eq}}\cos(k\theta)$
that produces the same $dS_y/dt$ as a misalignment of amplitude $A_q$.
The conversion factor
\begin{equation}
c_k \approx 6\,\mathrm{nT}\,\mu\mathrm{m}^{-1} \quad (k=1,\ldots,7,\;\pm17\%)
\label{eq:ck}
\end{equation}
is nearly $k$-independent: the spin couples to the ring-average impulse
rather than the local peak, so the FODO fill factor $f_\mathrm{fill}\sim3\%$
and gradient $G=\SI{0.21}{\tesla\per\meter}$ set the overall scale.
The local quad-centre field is $\sim\!34\times$ larger than $B_{x,\mathrm{eq}}$,
confirming that spin tracking is the appropriate reference.

Combining Eq.~(\ref{eq:ck}) with the BPM-offset sensitivity
$\sigma(A_k)\approx\sigma_b/\|M_k\|$ (Section~\ref{sec:whiteness}) gives a
measurement floor expressed directly in nT:
\begin{equation}
\sigma(B_{x,\mathrm{eq}}^{(k)}) \approx \frac{c_k\,\sigma_b}{\|M_k\|}.
\label{eq:bxsens}
\end{equation}
At $\sigma_b=\SI{100}{\micro\meter}$: $\sigma(B_{x,\mathrm{eq}}^{(2)})\approx
c_2\times\SI{0.6}{\micro\meter}\approx\SI{3.3}{\nano\tesla}$; at
$\sigma_b=\SI{20}{\micro\meter}$ (post-BBA): $\approx\SI{0.65}{\nano\tesla}$.
For $k=3$ the higher Fourier noise ($\|M_3\|=43.6$ versus $\|M_2\|=167$) demands
$\sigma_b\lesssim\SI{5}{\micro\meter}$ to reach $\SI{1}{\nano\tesla}$;
$\sigma_b\approx\SI{20}{\micro\meter}$ (post-BBA) gives
$\approx\SI{4}{\nano\tesla}$.

\subsection{Reconstruction performance and diagnostics}

Under $\sigma_b=\SI{100}{\micro\meter}$ the CLEAN reconstruction returns
\begin{equation}
A_{k=2} = 9.94\pm0.66~\si{\micro\meter}
\quad(\text{true }10),\qquad
\Delta\phi = \SI{0.046}{\radian},
\end{equation}
a $0.6\%$ amplitude error. Three diagnostics confirm the mechanism rather
than an artefact:
\begin{itemize}
\item \textbf{Signal vs.\ offset.}
$\|R\,\Delta q\|=\SI{5271}{\micro\meter}$,
$\|R\,\Delta q_{k=2}\|=\SI{1669}{\micro\meter}$, and
$\|\mathbf{b}\|=\SI{611}{\micro\meter}$: even the isolated $k=2$ orbit
exceeds the total offset norm.
\item \textbf{Pure-offset leakage.} With $\Delta q=0$ (offset only),
CLEAN attributes a spurious $k=2$ amplitude of just
$0.72\pm\SI{0.41}{\micro\meter}$, i.e.\ the genuine signal exceeds the
offset leakage by $\sim\times14$.
\item \textbf{Offset scaling.} Sweeping $\sigma_b$ shows the $k=2$ error
scaling linearly with the offset, as expected if $\mathbf{b}$ enters
through its (small) projection; even at a pessimistic
$\sigma_b=\SI{200}{\micro\meter}$ the $k=2$ error is
$\sim\SI{1.2}{\micro\meter}$, still within the \SI{10}{\micro\meter} target.
\end{itemize}

\subsection{Sensitivity to the response-matrix model}

In a real ring $R$ is known only approximately. Modelling a gradient
error $R_\mathrm{model}=R\,\mathrm{diag}(1+\varepsilon)$ with
$\varepsilon\sim\mathcal{N}(0,\sigma_\mathrm{model})$, the $k=2$ amplitude
estimate acquires a relative error that scales as
$\delta A_{k=2}/A_{k=2}\sim\sigma_\mathrm{model}\times\kappa_2$, where
$\kappa_2$ is the effective condition number of the $k=2$ column of $M$
under gradient perturbations. Numerically, the $k=2$ error
stays within the \SI{10}{\micro\meter} target for
$\sigma_\mathrm{model}=\delta K/K\lesssim3\text{--}4\%$, a tolerance
achievable with a beta-beat / LOCO-type calibration.
Figure~\ref{fig:model_error} quantifies the full error budget.
Three BPM-offset levels ($\sigma_b = 0$, $20$, $100\,\mu$m) are overlaid:
offset noise dominates at any realistic model accuracy, and model error
alone is negligible across the entire $0$--$10\%$ range.
The \SI{10}{\micro\meter} target is met for $\sigma_\mathrm{model}\lesssim10\%$
even at $\sigma_b=\SI{100}{\micro\meter}$,
confirming that the tolerance is driven by BPM offsets, not by imperfect
knowledge of $R$.

\subsection{Scope: targeted estimation, not harmonic discovery}

The method estimates the amplitude and phase of a \emph{prescribed} target
harmonic; it is not a harmonic-discovery tool. This is a deliberate scope
choice consistent with the operational objective stated in the Introduction:
the target $k=2$ is fixed in advance by the spin-tracking result of
Sec.~\ref{sec:falseedm}, not searched for. For completeness we note the
behaviour when the method is instead asked to discover content blindly:
because $\mathbf{b}$ projects a small amount onto every candidate mode,
CLEAN with the oracle-free candidate set $k=1,\dots,12$ flags all of them
above the \SI{1}{\micro\meter} level, so genuine harmonics ($k=2,4,6,8$)
cannot be distinguished from offset-induced spurious ones on amplitude
alone. Crucially, this does \emph{not} corrupt the targeted $k=2$
\emph{amplitude} (the modes are orthogonal under $R$); it only means the
method should not be repurposed as a general harmonic-detection scheme.
Such a discovery extension is discussed in Section~\ref{sec:future}.

\subsection{BPM-offset robustness without a whiteness assumption}
\label{sec:whiteness}

The Fourier-basis formulation replaces ``how large is $\mathbf{b}$?'' with
``how large is the projection of $\mathbf{b}$ onto the $k=2$ orbit pattern?''
For the fitting matrix $M=RF$, the contamination of the $k=2$ amplitude
estimate by \emph{any} deterministic offset is
\begin{equation}
\delta a_{k=2} = \frac{\mathbf{b}\cdot\hat{m}_{k=2}}{\|M_{k=2}\|}
               = \frac{\mathbf{b}\cdot\hat{m}_{k=2}}{167},
\label{eq:bias}
\end{equation}
where $\hat{m}_{k=2}=M_{k=2}/\|M_{k=2}\|$ is the unit vector along the
$k=2$ column of $M$. No assumption about the distribution of $\mathbf{b}$
enters Eq.~(\ref{eq:bias}). Strictly, this expression is exact only when the
columns of $M$ are mutually orthogonal; the general result is
$\delta\hat{a}=(M^TM)^{-1}M^T\mathbf{b}$.
Because $|\langle RF_2,RF_k\rangle|<0.01$ for all $k\neq2$ (Section~\ref{sec:results}),
Eq.~(\ref{eq:bias}) is an excellent approximation—the cross-correlation
correction to the $k=2$ estimate is below 1\%.

To stress-test this we set $\mathbf{b}=A\hat{m}_{k=2}$, the worst-case
fully $k=2$-orbit-aligned offset. Table~\ref{tab:whiteness} shows that
theory and simulation agree to four significant figures at all tested
amplitudes. An offset of $A=\SI{100}{\micro\meter}$ induces a
bias of \SI{0.60}{\micro\meter}, well within the \SI{10}{\micro\meter} target;
the safe limit $A<\SI{1669}{\micro\meter}$ is 16.7$\times$ above the
realistic value.

For white $\mathbf{b}$ with r.m.s.\ $\sigma_b$, the $k=2$ estimate standard
deviation is $\sigma_b/\|M_{k=2}\|=\sigma_b/167$; at the realistic
$\sigma_b=\SI{100}{\micro\meter}$ this gives $\lesssim\SI{0.6}{\micro\meter}$,
fully within target.

Three distinct amplitude scales, shown in Fig.~\ref{fig:bpm_fourier},
quantify the robustness end-to-end:
\begin{enumerate}
\item \textbf{Signal orbit ($\SI{1670}{\micro\meter}$).}
  A $\SI{10}{\micro\meter}$ $k=2$ misalignment passes through $R$ and is
  resonantly amplified: $A\|M_{k=2}\|=10\times167=\SI{1670}{\micro\meter}$
  in BPM space.
\item \textbf{BPM-offset Fourier level ($\approx\SI{26}{\micro\meter}$).}
  Because $\mathbf{b}$ bypasses $R$ entirely and is white, its projection
  onto any FODO-antisymmetric mode $F_k$ has r.m.s.\
  $\sigma_b\sqrt{\pi/48}\approx\SI{26}{\micro\meter}$---distributed equally
  across all $k$, with no peak at $k=2$.
  The signal orbit ($\SI{1670}{\micro\meter}$) is already $64\times$ larger than
  this raw Fourier level.
\item \textbf{Estimator contamination floor ($\approx\SI{0.6}{\micro\meter}$).}
  The estimator projects $\mathbf{b}$ onto the $k=2$ column of $M=RF$
  (not onto $F_{k=2}$ directly), dividing by $\|M_{k=2}\|=167$ a second
  time: $\sigma_b/\|M_{k=2}\|=\SI{100}{\micro\meter}/167\approx\SI{0.6}{\micro\meter}$.
\end{enumerate}
The hierarchy
$\SI{1670}{\micro\meter}\gg\SI{26}{\micro\meter}\gg\SI{0.6}{\micro\meter}$
shows that the orbit amplification of $\|M_{k=2}\|$ works doubly in
the estimator's favour: it enlarges the genuine signal in BPM space
while simultaneously suppressing the offset contribution in the
estimate.

The key intuition is that \emph{the estimator does not measure how large the
orbit is---it measures how closely the orbit resembles the $k=2$ fingerprint.}
A genuine $k=2$ signal shares exactly that fingerprint, so the element-wise
product $y_j\hat{m}_{k=2,j}$ is everywhere positive and the sum is large.
A random BPM offset carries no such structure: its products with
$\hat{m}_{k=2}$ alternate in sign and nearly cancel.
Figure~\ref{fig:matched_filter} illustrates this directly for a
$\sigma_b=\SI{100}{\micro\meter}$ offset alongside a $\SI{10}{\micro\meter}$
$k=2$ signal, recovering $\hat{a}_{k=2}=10.0\,\mu$m and $0.4\,\mu$m respectively.
Figure~\ref{fig:whiteness_spec} confirms that the recovered-amplitude spectrum
of a white offset is flat and broadband, with no spurious peak, and the floor
is lowest precisely at $k=2$ ($\SI{0.7}{\micro\meter}$) because $\|M_{k=2}\|=167$
is the largest column norm.

\begin{table}[t]
\centering
\caption{Worst-case bias test: $\mathbf{b}=A\hat{m}_{k=2}$
(offset fully aligned with the $k=2$ orbit pattern).
Theory: $\delta a=A/167$; CLEAN and lstsq reproduce theory to
four significant figures.}
\label{tab:whiteness}
\begin{tabular}{cccc}
\toprule
$A$ (\si{\micro\meter}) & Theory $\delta a$ (\si{\micro\meter}) &
CLEAN (\si{\micro\meter}) & lstsq (\si{\micro\meter}) \\
\midrule
  20 & 0.12 & $+$0.12 & $+$0.12 \\
  50 & 0.30 & $+$0.30 & $+$0.30 \\
 100 & 0.60 & $+$0.60 & $+$0.60 \\
 200 & 1.20 & $+$1.20 & $+$1.20 \\
1669 & 10.0 & $+$10.0 & $+$10.0 \\
\bottomrule
\end{tabular}
\end{table}

%==============================================================
\section{Discussion and future work}
\label{sec:future}

The single-orbit, response-matrix reconstruction trades the offset-cancelling
guarantee of kmod for a structural separation that is effective precisely
for the low-order, tune-amplified harmonic that matters for the false EDM.
It is important to recall the framing stated in the Introduction: this is
not a general-purpose alignment tool but a targeted estimator for a single
physical observable---the $k=2$ amplitude that drives the false EDM.
A note on kmod: with \emph{uniform} modulation of all 48 quadrupoles
simultaneously, $\Delta R=\varepsilon R$ and $\kappa(\Delta R)=\kappa(R)\approx249$,
so in principle full-ring kmod achieves equivalent conditioning without the
BPM-offset problem. In practice, 48 simultaneous precision gradient steps
are operationally demanding; the single-orbit method requires only one
standard orbit measurement and no active gradient manipulation.
The following items remain open.

\begin{itemize}
  \item Validated at the reference offset $\sigma_b=\SI{100}{\micro\meter}$
    via the worst-case structured test (Section~\ref{sec:whiteness}): the $k=2$ bias
    remains $\lesssim\SI{0.6}{\micro\meter}$ with no whiteness assumption.
    \todo{Include non-zero readout noise $\boldsymbol{\eta}$ in the full error budget.}
\item \todo{Discriminate genuine from offset-induced harmonics without
oracle knowledge---e.g.\ calibrate the spurious-amplitude threshold from
the pure-offset diagnostic, or impose a low-$k$ physical prior.}
\item \todo{Multi-orbit extension: combine orbits taken at different times
so that a slowly drifting $\mathbf{b}$ averages down while $\Delta q$ stays
fixed, further suppressing the offset.}
\item \todo{Response-matrix calibration on a real lattice: quantify how
well $R$ can be measured (LOCO/beta-based BBA) and propagate the residual
$\delta K/K$ to the $k=2$ error budget.}
\item \todo{Connect $\Delta q_{k=2}$ to the false-EDM functional
$\Phi_\mathrm{spin}\propto\sum_j K_jL_j(y_\mathrm{CO}(s_j)-\Delta y_j)$,
i.e.\ propagate the measured harmonic to a spin-level systematic and
define the closure tolerance.}
\item \todo{Extend to the horizontal plane and to skew coupling (quad
tilt), reusing the horizontal response matrix $R_x$ already built.}
\end{itemize}

%==============================================================
\section{Conclusion}
\label{sec:conclusion}

We have shown, in simulation, that the $k=2$ harmonic of the
quadrupole-misalignment pattern---which we first establish by spin tracking
to be the dominant driver of the false EDM in a pEDM frozen-spin ring---can
be estimated to $0.6\%$ in amplitude and \SI{0.046}{\radian} in phase from a
single closed-orbit measurement in the presence of a \SI{100}{\micro\meter}
BPM offset, without gradient modulation. The broader message is operational:
when the objective is suppression of the dominant false-EDM harmonic rather
than metrological alignment, a full reconstruction of the misalignment
vector is unnecessary---it suffices to measure the one mode that matters and
then null it with dedicated correctors (a correction scheme left to future
work). The result rests on three computable properties of the
response matrix: tune-proximity amplification of $k=2$ (column norm
$\|M_{k=2}\|=167$ vs.\ 16 for $k=4$), the consequent inversion of the
misalignment-amplitude hierarchy at the orbit level, and near-orthogonality
of the harmonic orbit responses that prevents contaminant leakage.
Critically, the method makes no assumption about the statistical structure
of the BPM offset: even if the full \SI{100}{\micro\meter} offset is
entirely aligned with the $k=2$ orbit pattern (worst case), the induced
bias is only \SI{0.60}{\micro\meter}; the idealized worst-case limit is $\SI{1669}{\micro\meter}$,
16.7$\times$ above the realistic value (Section~\ref{sec:whiteness}).
We also documented and corrected a $10^3$ unit-scaling error in the
response-matrix pipeline. Within the targeted-estimation scope adopted here
the method is complete; the natural extensions---propagating the measured
harmonic to a spin-level error budget, and a discovery mode that
discriminates genuine from offset-induced harmonics without prior
knowledge---are set out in Section~\ref{sec:future}.

%==============================================================
\section*{Acknowledgements}
\todo{acknowledgements / funding}

\begin{thebibliography}{9}
\bibitem{omarov} Z.~Omarov, H.~Davoudiasl, S.~Hac\i\"omero\u{g}lu,
V.~Lebedev, W.~M.~Morse, Y.~K.~Semertzidis, A.~J.~Silenko,
E.~J.~Stephenson, and R.~Suleiman, \emph{Comprehensive symmetric-hybrid
ring design for a proton EDM experiment at below $10^{-29}\,e\cdot$cm},
Phys.\ Rev.\ D \textbf{105} (2022) 032001.
\bibitem{clean} J.~A.~H\"ogbom, \emph{Aperture Synthesis with a
Non-Regular Distribution of Interferometer Baselines}, Astron.\ Astrophys.\
Suppl.\ \textbf{15} (1974) 417.
\bibitem{loco} J.~Safranek, \emph{Experimental determination of storage
ring optics using orbit response measurements}, Nucl.\ Instrum.\ Meth.\ A
\textbf{388} (1997) 27.
\end{thebibliography}

\end{document}
